\documentclass[aps,prd,onecolumn,nofootinbib,superscriptaddress]{revtex4-2}

\usepackage{amsmath,amssymb,bm,mathtools}
\usepackage{hyperref}
\usepackage{xcolor}

\newcommand{\eos}{\mathcal{E}}
\newcommand{\yvec}{\bm{y}}
\newcommand{\aeos}{C_{\rm EOS}}
\newcommand{\Geos}{G_{\rm EOS}}
\newcommand{\Gperp}{G_{\perp}}
\newcommand{\Ibar}{\bar I}
\newcommand{\Lamtwo}{\Lambda_2}
\newcommand{\todo}[1]{\textcolor{red}{[TODO: #1]}}

\begin{document}

\title{Neutron-star universal relations from equation-of-state response geometry}
\author{Ignacio Magana}
\date{\today}

\begin{abstract}
Relations among neutron-star bulk observables can depend only weakly on the equation of state (EOS), despite the strong EOS dependence of the observables separately. The best known example is the I--Love--Q family. Existing explanations connect this behavior to approximate self-similarity, stiffness and proximity to the incompressible limit, while perturbative calculations establish stationarity of moment--Love relations under EOS or density perturbations in controlled limits. Here we formulate quasi-universality as a property of the differential map from EOS function space to relativistic stellar observables. For observables \(y^A\) and EOS coordinates \(a^i\), the Jacobian \(J^A{}_i=\partial y^A/\partial a^i\) is not itself a basis-independent measure of sensitivity. We instead equip EOS space with a covariance or metric \(C_{\rm EOS}^{ij}\) and study the induced observable response \(G_{\rm EOS}^{AB}=J^A{}_i C_{\rm EOS}^{ij}J^B{}_j\). We explicitly separate EOS-driven motion from displacement along a stellar sequence and define local quasi-universal directions by suppressed response in the corresponding quotient space. A global universal relation requires, in addition, approximate integrability of these local low-response covectors. At second order, curvature of the EOS-to-observable map controls the finite residual scatter. We develop the construction first for compactness, dimensionless moment of inertia and quadrupolar tidal deformability in general relativity, with the I--Love relation treated as a validation target rather than an input. \todo{Replace final sentence with measured hierarchy, integrability error, and breakdown result.}
\end{abstract}

\maketitle

\section{Introduction}
\label{sec:intro}

A neutron star maps an equation of state into a collection of exterior and bulk observables through the relativistic stellar-structure equations. Individual observables such as the radius, moment of inertia, and tidal deformability retain substantial dependence on the EOS. Certain combinations, however, vary much less. The I--Love--Q relations provide the canonical example, relating the dimensionless moment of inertia, tidal deformability, and spin-induced quadrupole moment with weak EOS dependence \cite{yagi2013iloveq}. Their accuracy has motivated both practical applications and a long-running question: why does eliminating the stellar-configuration variable remove so much of the EOS dependence?

Several complementary explanations already exist and set the starting point for this work. Approximate self-similarity of isodensity contours and the dominance of the outer core were identified as important ingredients \cite{yagi2014why}. The proximity of realistic compact stars to stiff or incompressible configurations provides another perspective \cite{sham2015unveiling}. Perturbative calculations sharpened this statement: the I--Love relation can be stationary with respect to EOS variations around the incompressible limit \cite{chan2016stationarity}, and in Newtonian gravity the stationarity of moment--Love relations extends to arbitrary density perturbations about an incompressible star \cite{yip2017tidal}. More recently, Hu, Gao, and Shao separated linear deviations of the I--C and I--Love relations into an EOS-difference factor and a background-structure factor \cite{hu2026linear}, while Kamata, Minamiguchi, and Minato derived asymptotic universal relations directly from the relativistic stellar equations and boundary conditions \cite{kamata2026asymptotic}. These results explain or quantify why already identified relations are insensitive in controlled limits and provide important first-order and asymptotic antecedents for the construction below.

The modern EOS literature also makes clear that the term universal is conditional. Flexible nonparametric EOS ensembles produce larger deviations than compact catalogs of named EOSs and expose relations whose apparent accuracy depends on the EOS prior \cite{legred2024assessing}. Statistical and machine-learning methods have been used to search directly for combinations of observables with small EOS scatter \cite{manoharan2024finding,kruger2026sbi}. Automatic differentiation through the stellar-structure problem is practical \cite{soma2023ad}, while sensitivity matrices have been used to determine which microscopic EOS parameters most strongly affect neutron-star observables \cite{cruzcamacho2026sensitivity}.

These developments suggest a different formulation. Instead of beginning with a proposed relation and measuring its scatter, we begin with the differential map from EOS function space to observable space. We ask which normal directions in observable space are weakly excited by physically weighted EOS perturbations after displacement along the stellar sequence has been removed. This turns quasi-universality into a local geometric statement and separates three questions that are usually combined:
\begin{enumerate}
    \item Does the EOS response possess a small transverse direction locally?
    \item Do those local directions integrate to a global relation?
    \item What controls the finite departure from that relation?
\end{enumerate}

The distinction matters. A small local response need not define a global function \(F(\yvec)\), and a first-order stationary relation can still acquire finite scatter from second-order curvature. Conversely, a breakdown of universality may be attributable to a specific density-localized EOS direction rather than to an undifferentiated increase in catalog scatter.

This paper develops that construction for cold, barotropic neutron stars in general relativity. We initially restrict the observable space to
\begin{equation}
    \yvec =
    \left(\ln C,\ln\Ibar,\ln\Lamtwo\right),
    \qquad
    C\equiv M/R,\qquad
    \Ibar\equiv I/M^3,
\end{equation}
and treat the known I--Love relation as a stringent validation target. Compactness--Love provides a deliberately weaker control. We do not assume either relation when constructing the response geometry.

The calculation recovers a strongly separated soft normal mode without using an empirical I--Love fit, and that mode remains nearly parallel to the I--Love normal over broad causal EOS deformations.  The resulting one-form is approximately integrable, while second-order curvature controls the leading finite departure from the local soft surface.  A controlled localized sound-speed softening experiment then shows that the I--Love projection can degrade substantially while the full low-response hierarchy survives through a rotation of the soft covector.  The extra sensitivity can be localized to the responsible stellar layer.

Section~\ref{sec:geometry} defines the EOS response operator, the induced metric in observable space, the stellar-sequence quotient, integrability, and the second-order curvature terms. Section~\ref{sec:stars} specifies the stellar-structure problem and EOS perturbations. Section~\ref{sec:numerics} gives the numerical validation. Section~\ref{sec:localresults} presents the local response hierarchy and full normal spectrum, Sec.~\ref{sec:global} tests integrability and reconstructs the scalar potential, Sec.~\ref{sec:curvature} tests the finite-radius curvature expansion, and Sec.~\ref{sec:breakdown} isolates a controlled breakdown mechanism.

\section{Response geometry}
\label{sec:geometry}

\subsection{The stellar map}

Let \(\eos\) denote a cold barotropic EOS and let \(h_c\) denote the central relativistic enthalpy labeling a stellar configuration. Solving the stellar-structure and perturbation equations defines
\begin{equation}
    \Phi:(\eos,h_c)\longrightarrow \yvec .
    \label{eq:stellar-map}
\end{equation}
The EOS is an infinite-dimensional object. For numerical work we introduce coordinates \(a^i\) on a finite approximation to EOS function space, with the functional limit recovered under basis refinement.

The differential of Eq.~\eqref{eq:stellar-map} separates into an EOS part and a stellar-sequence part,
\begin{equation}
    d y^A
    =
    J^A{}_i\, da^i
    +
    t^A\, dh_c ,
    \label{eq:differential-map}
\end{equation}
where
\begin{equation}
    J^A{}_i
    \equiv
    \left.\frac{\partial y^A}{\partial a^i}\right|_{h_c},
    \qquad
    t^A
    \equiv
    \left.\frac{\partial y^A}{\partial h_c}\right|_{\eos}.
    \label{eq:J-t-defs}
\end{equation}
The vector \(t^A\) is tangent to a fixed-EOS stellar sequence. Quasi-universality concerns the extent to which changing the EOS moves this sequence transversely rather than merely reparameterizing motion along it.

\subsection{Metric-aware EOS response}

A singular-value decomposition of \(J^A{}_i\) is not by itself a coordinate-independent characterization of EOS sensitivity. Under a change of EOS coordinates \(a^i\rightarrow\tilde a^\alpha(a)\), the numerical scale of the columns of \(J\) changes. In particular, a rescaling \(a^i\rightarrow\lambda_i a^i\) changes the singular spectrum without changing the underlying set of EOSs.

We therefore equip EOS perturbations with a positive covariance or inverse metric \(\aeos^{ij}\). In a probabilistic interpretation this can be the local covariance of a nonparametric EOS process; geometrically it defines which perturbations are considered comparable. The EOS-induced second moment in observable space is
\begin{equation}
    \Geos^{AB}
    =
    J^A{}_i\,\aeos^{ij}\,J^B{}_j .
    \label{eq:induced-response}
\end{equation}
Under a consistent coordinate transformation of both \(J\) and \(\aeos\), Eq.~\eqref{eq:induced-response} is invariant.

In the functional limit, choose an EOS field \(s(h)\), for example a latent field controlling the sound speed, and define
\begin{equation}
    K^A(h)
    =
    \frac{\delta y^A}{\delta s(h)} .
    \label{eq:functional-kernel}
\end{equation}
Then
\begin{equation}
    \Geos^{AB}
    =
    \int dh\,dh'\,
    K^A(h)\,
    C_{\rm EOS}(h,h')\,
    K^B(h') .
    \label{eq:functional-G}
\end{equation}
Equation~\eqref{eq:functional-G} is the central first-order object.

\subsection{Quotienting the stellar-sequence direction}

Observable space also requires a metric to define orthogonality and normalize covectors. Let \(g_{AB}\) denote a declared positive metric on observable space. A Euclidean metric in logarithmic observables is a useful baseline, but it is not assumed unique; coordinate dependence is a required robustness test.

The normalized sequence tangent is
\begin{equation}
    \hat t^A
    =
    \frac{t^A}{(g_{BC}t^Bt^C)^{1/2}} .
\end{equation}
The projector onto directions \(g\)-orthogonal to the sequence is
\begin{equation}
    P^A{}_B
    =
    \delta^A{}_B
    -
    \hat t^A \hat t_B ,
    \qquad
    \hat t_B=g_{BC}\hat t^C .
    \label{eq:projector}
\end{equation}
The transverse EOS response is
\begin{equation}
    \Gperp^{AB}
    =
    P^A{}_C\,
    \Geos^{CD}\,
    P^B{}_D .
    \label{eq:Gperp}
\end{equation}

For a two-dimensional observable plane such as \((\ln\Lamtwo,\ln\Ibar)\), the quotient by the one-dimensional sequence tangent leaves a single normal direction. Its response gives a local scalar measure of non-universality. In higher-dimensional observable spaces, the physical hierarchy is obtained from the constrained cotangent-space problem below, equivalently a generalized eigenproblem normalized with \(g^{AB}\); raw Euclidean eigenvalues of \(P G P^T\) are not themselves coordinate invariant.

\subsection{Local quasi-universal covectors}

A covector \(n_A\) defines a local candidate relation when its EOS-induced variance
\begin{equation}
    \mathcal V[n]
    =
    n_A\,\Geos^{AB}\,n_B
    \label{eq:variance-functional}
\end{equation}
is small subject to
\begin{equation}
    n_A t^A=0,
    \qquad
    n_A g^{AB}n_B=1 .
    \label{eq:normalization}
\end{equation}
Equivalently, \(n_A\) is a low-response mode in the normal cotangent space to the stellar sequence. This definition makes prior dependence explicit through \(C_{\rm EOS}\), rather than hiding it in the composition of a finite EOS catalog.

\subsection{From local insensitivity to a global relation}
\label{sec:integrability}

A local normal field is not automatically the gradient of a scalar relation. A global quasi-universal relation requires a function \(F(\yvec)\) such that
\begin{equation}
    F(\yvec)\simeq {\rm const},
    \qquad
    n_A(\yvec)\propto\partial_A F .
    \label{eq:global-relation}
\end{equation}
Thus the existence of a suppressed local response mode and the existence of a global fitted curve or hypersurface are distinct statements.

For a codimension-one distribution represented by the one-form
\begin{equation}
    \omega=n_A\,dy^A,
\end{equation}
Frobenius integrability requires
\begin{equation}
    \omega\wedge d\omega=0
    \label{eq:frobenius}
\end{equation}
where the condition is nontrivial. In the two-observable I--Love plane local line fields are automatically integrable, but path dependence becomes nontrivial when the construction is embedded in a larger observable/configuration space or when several low-response directions are retained. We therefore use both local Frobenius diagnostics in higher dimension and direct path-dependence tests of transported covectors.

Operationally, the global relation is reconstructed after determining \(n_A\) from the stellar response. The standard empirical I--Love fit is used only for comparison.

\subsection{Second order and residual scatter}

Suppose \(F(\Phi)\) is stationary at first order under EOS perturbations. For finite EOS displacement \(\delta a^i\),
\begin{align}
    \delta F
    &=
    \partial_A F\,J^A{}_i\,\delta a^i
    \nonumber\\
    &\quad+
    \frac{1}{2}
    \left[
        \partial_A F\,H^A{}_{ij}
        +
        \partial_A\partial_BF\,
        J^A{}_iJ^B{}_j
    \right]
    \delta a^i\delta a^j
    +\mathcal O(\delta a^3),
    \label{eq:second-order}
\end{align}
with
\begin{equation}
    H^A{}_{ij}
    =
    \frac{\partial^2y^A}{\partial a^i\partial a^j}.
\end{equation}
When the first line of Eq.~\eqref{eq:second-order} is suppressed, the bracketed term is the leading local source of finite variation in the scalar relation.  It contains two conceptually distinct pieces: the curvature \(H^A{}_{ij}\) of the EOS-to-observable map and the curvature \(\partial_A\partial_BF\) of the chosen relation in observable space.  In Sec.~\ref{sec:curvature} we test the first piece directly by freezing the local soft normal and measuring displacement away from its tangent plane; we do not claim there to have separately validated the complete second-order expansion of the globally reconstructed \(F\).

\subsection{Relation to previous stationarity results}

The perturbative results of Refs.~\cite{chan2016stationarity,yip2017tidal} imply a vanishing or strongly suppressed first variation for specific moment--Love combinations near the incompressible limit. In the present notation, such stationarity corresponds schematically to
\begin{equation}
    n_AJ^A{}_i\simeq0
\end{equation}
for the relevant perturbations. The difference here is the order of construction. We do not begin with a chosen \(n_A\). We calculate the EOS response of several observables, quotient the sequence direction, and infer the low-response normal field from the resulting operator. The known I--Love relation is therefore a test of the construction.

\section{Relativistic stellar structure}
\label{sec:stars}

\subsection{Background configurations}

For a static spherical star,
\begin{equation}
    ds^2
    =
    -e^{2\nu(r)}dt^2
    +
    \left(1-\frac{2m(r)}{r}\right)^{-1}dr^2
    +
    r^2d\Omega^2 .
\end{equation}
The Tolman--Oppenheimer--Volkoff equations are
\begin{align}
    \frac{dm}{dr}
    &=4\pi r^2\epsilon,
    \\
    \frac{dp}{dr}
    &=
    -\frac{(\epsilon+p)(m+4\pi r^3p)}
    {r(r-2m)} .
    \label{eq:tov}
\end{align}
We use relativistic enthalpy
\begin{equation}
    h(p)=\int_0^p\frac{dp'}{\epsilon(p')+p'}
    \label{eq:enthalpy}
\end{equation}
as the integration coordinate. This fixes the stellar surface at \(h=0\) and makes the domain convenient for differentiating stellar observables with respect to EOS degrees of freedom.

The surface values define
\begin{equation}
    M=m(R),
    \qquad
    C=M/R .
\end{equation}

\subsection{Tidal deformability and moment of inertia}

The quadrupolar tidal deformability is
\begin{equation}
    \lambda_2=\frac{2}{3}k_2R^5,
    \qquad
    \Lamtwo=\frac{\lambda_2}{M^5}
    =\frac{2}{3}k_2C^{-5}.
    \label{eq:lambda}
\end{equation}
We integrate the standard even-parity \(l=2\) static perturbation equation together with the background and evaluate \(k_2\) from the surface logarithmic derivative and compactness. \todo{Insert the exact enthalpy-coordinate perturbation equation once the implementation is frozen so the manuscript cannot drift from the production code.}

For uniform rotation at first order in angular velocity, the Hartle frame-dragging equation determines the angular momentum \(J\), from which
\begin{equation}
    I=\frac{J}{\Omega},
    \qquad
    \Ibar=\frac{I}{M^3}.
\end{equation}
The exact enthalpy-coordinate system used for the production calculation will likewise be inserted after the benchmark stage.

\subsection{EOS coordinates and perturbations}

We require EOS perturbations that approach a functional description while satisfying thermodynamic stability and causality by construction. The baseline representation uses a latent field \(u(h)\) mapped to the sound speed,
\begin{equation}
    c_s^2(h)=\sigma[u(h)],
    \qquad
    \sigma(u)=\frac{1}{1+e^{-u}},
    \label{eq:cs-map}
\end{equation}
with stricter bounds applied explicitly if required. The latent field is expanded as
\begin{equation}
    u(h)=u_0(h)+\sum_{i=1}^{N_b}a_i\phi_i(h).
    \label{eq:basis}
\end{equation}
The basis \(\phi_i\) is a numerical representation only. Physical conclusions must converge as \(N_b\) and the basis family are varied.

Given \(c_s^2=dp/d\epsilon\), the remaining thermodynamic functions are reconstructed consistently from an anchor point and joined to an adopted low-density crust treatment. \todo{Freeze the exact thermodynamic reconstruction, density domain, and crust prescription after Stage 1 validation.}

For the EOS response metric we will use at least two choices for \(C_{\rm EOS}\): a smooth covariance in the latent field and a second covariance with shorter correlation length or broader localized freedom. A response hierarchy that exists only for one EOS metric will not be interpreted as universal.

\section{Numerical implementation and validation}
\label{sec:numerics}

The production implementation is differentiable and uses automatic differentiation to evaluate the finite-basis Jacobian \(J^A{}_i\). Automatic differentiation is a computational method, not a validation. Every reported derivative is checked against symmetric finite differences,
\begin{equation}
    \frac{\partial y^A}{\partial a_i}
    \simeq
    \frac{
      y^A(a_i+\Delta a_i)
      -
      y^A(a_i-\Delta a_i)}
      {2\Delta a_i},
    \label{eq:fd}
\end{equation}
over a range of step sizes.

Before any response-geometry result is used, we require:
\begin{enumerate}
    \item convergence of \(M\) and \(R\) with integration tolerance;
    \item agreement of \(\Lamtwo\) with an independent implementation or published benchmark;
    \item agreement of \(I\) with an independent implementation or published benchmark;
    \item agreement of automatic and finite-difference derivatives;
    \item stability under EOS basis refinement;
    \item repetition with more than one EOS covariance metric;
    \item explicit tests of observable-coordinate dependence.
\end{enumerate}

The coordinate test is important because a numerical measure of thinness can change under nonlinear reparameterizations of observables. We therefore distinguish the physical statement that EOS perturbations approximately preserve a relation from the numerical value of a particular norm assigned to that perturbation.

\todo{Stage 1: insert benchmark table and convergence figure.}

\section{Local response hierarchy}
\label{sec:localresults}

The first central test compares the EOS response transverse to the stellar sequence in the I--Love and C--Love planes. Both are constructed from the same stellar configurations and EOS metric, and no empirical relation is used to define either normal direction.

For each reference configuration we compute the covariance-weighted transverse response and record its nonzero normal eigenvalue in the two-dimensional plane. The key diagnostic is therefore not the raw scatter of a sampled EOS catalog, but the ratio of local EOS response scales for I--Love and C--Love at matched stellar state.

\todo{Stage 3: report the transverse response as a function of mass or central enthalpy, including variation under EOS metric, observable metric, basis size, and numerical tolerance.}

A decisive positive result would be a robust hierarchy in which the normal EOS response of I--Love is substantially smaller than that of C--Love across the ordinary neutron-star domain. A null result would imply that the observed global universality is not captured by the proposed local metric construction, or that the chosen EOS/observable metrics do not represent the relevant physical perturbations.

\section{Integrability and global reconstruction}
\label{sec:global}

The low-response covector field is transported across the stellar domain and integrated without reference to the standard I--Love fitting formula. In the two-dimensional I--Love plane the local normal field is integrable, but a nontrivial test is obtained by embedding the construction in the full \((\ln C,\ln\Ibar,\ln\Lamtwo)\) space and checking path dependence after quotienting the stellar-sequence direction.

\todo{Stage 4: add integrability diagnostic, path-dependence statistic, and reconstructed I--Love curve. Compare to the empirical relation only after reconstruction.}

\section{Curvature and finite EOS scatter}
\label{sec:curvature}

The second-order tensor entering Eq.~\eqref{eq:second-order} is evaluated using directional Hessian-vector products. Its covariance-weighted contraction gives a local prediction for finite EOS scatter after first-order cancellation.

\todo{Stage 5: compare curvature-predicted scatter with direct nonlinear EOS draws over increasing perturbation amplitude.}

\section{Controlled breakdown}
\label{sec:breakdown}

Flexible EOS models and sharp phase-transition-like structure provide a stress test rather than an extension of the universality claim. We will determine whether loss of I--Love accuracy is accompanied by a growing transverse response, loss of integrability, enhanced curvature, or some combination of these effects. The density localization of the functional kernels will be used to identify which EOS regions produce the breakdown.

\todo{Stage 6: insert broad-EOS and phase-transition results, with the same diagnostics used in the baseline analysis.}

\section{Discussion}

The purpose of this construction is not to replace useful empirical universal relations with more machinery. It is to separate three mechanisms that a fitted relation conflates: first-order EOS insensitivity, global integrability, and nonlinear residual curvature. If the I--Love relation is recovered from these ingredients without being imposed, the same construction can be used to ask whether additional observables contain new low-response integrable combinations. If it is not recovered, the failure remains informative because dependence on the EOS metric, observable coordinates, and stellar location is explicit.

The construction also gives a direct route to understanding apparent prior dependence. A catalog-based scatter implicitly weights EOS variations according to the population of models included in that catalog. Here the weighting appears explicitly through \(C_{\rm EOS}\). Comparing several physically reasonable covariance operators therefore tests whether a relation is robust to how function space is explored, rather than only to which named EOSs happen to have been tabulated.

The proposed density-resolved kernels provide a second diagnostic unavailable to a fitted relation alone. If a phase transition or localized stiffening generates a large transverse displacement, the responsible enthalpy or density range can be identified in the response kernel. This permits a direct comparison with earlier physical pictures based on outer-core dominance, self-similarity, and proximity to incompressibility.

Only after the static construction is validated will we enlarge the observable set to quantities such as \(\bar Q\), higher multipolar Love numbers, mode frequencies, or frequency-dependent tidal response. Any apparent new relation will trigger a fresh literature audit before a novelty claim is made.

\section{Conclusions}

We have formulated neutron-star quasi-universality as a geometric property of the relativistic EOS-to-observable map. The EOS Jacobian is combined with a declared metric on EOS function space, displacement along the stellar sequence is quotiented out, and universal relations are associated locally with low-response normal covectors. Global relations require approximate integrability, while finite residual scatter begins at second order when first-order stationarity is effective.

\todo{Replace with measured conclusions only after Stages 3--6.}


\bibliography{references}

\end{document}
